\subsection{Representations of stellar algebras}

Let A be a stellar algebra. Denote by \(A_{+}\) the closed convex cone of positive elements of A (def. 6 of I, p. 115). A linear form \(\lambda\) on A is positive if and only if \(\lambda(\mathrm{A}_{+}) \subset \mathbf{R}_{+}\) (I, p. 118, th. 2).

PROPOSITION 7. --- Let A be a stellar algebra. Every positive linear form \(\lambda\) on A is continuous.

Let us first show that \(\lambda\) is bounded on the intersection of \(\mathrm{A}_{+}\) and the unit ball of A. Suppose this is not the case. Then there exists a sequence \((x_{n})_{n\geqslant 1}\) in \(\mathrm{A}_{+}\) such that \(\| x_n\| \leqslant 1\) and \(\lambda (x_n)\geqslant n\) for every integer \(n\geqslant 1\) . The series with general term \(n^{-2}x_{n}\) converges to an element \(x\) of A (cf. TG, IV, p. 33, example 3). For every integer \(\mathrm{N}\geqslant 1\) , since

\[
\sum_ {n = \mathrm{N} + 1} ^ {+ \infty} \frac {1}{n ^ {2}} x _ {n} \in \mathrm{A} _ {+},
\]

(I, p. 116, prop. 14) it follows that

\[
\sum_ {n = 1} ^ {\mathrm{N}} \frac {1}{n} \leqslant \sum_ {n = 1} ^ {\mathrm{N}} \frac {1}{n ^ {2}} \lambda (x _ {n}) = \lambda (x) - \lambda \left(\sum_ {n = \mathrm{N} + 1} ^ {+ \infty} \frac {1}{n ^ {2}} x _ {n}\right) \leqslant \lambda (x),
\]

which is absurd (TG, IV, p. 33, example 4).

Since every element of the unit ball of A is a linear combination with coefficients bounded by 1 of at most four positive elements of A of norm \(\leqslant 1\) (I, p. 96, lemma 2 and formula (4) of I, p. 117), we conclude that \(\lambda\) is continuous.

There exist involutive Banach algebras which admit positive linear forms that are not continuous (exercise 3, a) of V, p. 493).

PROPOSITION 8. --- Let A be a stellar algebra and let a be a nonzero element of A. There exists a positive linear form \(\lambda \in \mathrm{A}_{+}^{\prime}\) such that \(\lambda(a^{*}a) > 0\) .

Consider the real Banach space \(A_{h}\) of Hermitian elements of A. The set \(A_{+}\) is a closed convex pointed salient cone in \(A_{h}\) (prop. 14 of I, p. 116). The element \(a^{*}a\) is positive (th. 2 of I, p. 118) and nonzero, hence the Hermitian element \(-a^{*}a\) is not positive. By EVT, II, p. 42, cor. 5, there exists a continuous real linear form \(\lambda\in A_{h}^{\prime}\) such that \(\lambda(-a^{*}a)<0\) and \(\lambda(\mathrm{A}_{+})\subset\mathbf{R}_{+}\) . The linear form \(\lambda\) extends to a Hermitian C-linear form on A (cf. I, p. 98), which is positive and possesses the required property.

THEOREM 2 (Gelfand--Naimark). --- Let A be a stellar algebra. There exists a Hilbert space E and an isometric morphism of involutive algebras \(\varphi\) from A into \(\mathcal{L}(\mathrm{E})\) . If A is unital, there exists such a unital morphism.

For every \(b \in \mathrm{A} - \{0\}\) , let \(\lambda_b\) be a continuous positive linear form on A such that \(\lambda_b(b^*b) > 0\) (prop. 8). Since A admits an approximate unit (I, p. 121, prop. 18), there exists a Hilbertian realization \((\mathrm{E}_b, x_b, \varphi_b)\) of \(\lambda_b\) (prop. 5). If A is unital, one may suppose that \(\varphi_b\) is unital (loc. cit.). We have \(\| \varphi_b(b) \|^2 = \lambda_b(b^*b) \neq 0\) .

Let E be the external Hilbert sum of the spaces \(E_{b}\) for \(b\) belonging to \(A - \{0\}\) . For every \(a \in A\) , denote by \(\varphi(a) \in \mathcal{L}(E)\) the unique continuous linear map whose restriction to \(E_{b}\) coincides with \(\varphi_{b}(a)\) for all \(b \in A - \{0\}\) . The map \(a \mapsto \varphi(a)\) is a morphism of involutive algebras; it is injective, hence isometric (prop. 9 of I, p. 112), and satisfies \(\varphi(1) = 1_{\mathrm{E}}\) if A is unital. This concludes the proof.

*The category whose objects are the stellar algebras and whose morphisms are the morphisms of involutive algebras is therefore equivalent to the category whose objects are the closed involutive subalgebras of the endomorphism algebras of Hilbert spaces, and whose morphisms are the morphisms of involutive algebras.*

